\section{Related works}
\subsection{Diffusion model}

Diffusion models have recently become a practical choice for modeling complex, high-dimensional distributions. Early denoising diffusion probabilistic models formulate generation as the reverse of a gradual noising process, where a neural network learns to predict the noise added at each diffusion step~\cite{ddpm}. Score-based generative models provide a closely related view by learning the score of the data distribution and sampling through stochastic differential equations~\cite{score}. These models were first popularized in image generation, but the same denoising formulation is also useful when the generated object is not an image but a sequence of robot actions.

Diffusion Policy~\cite{dp} applies this idea to visuomotor control. Instead of predicting a single deterministic action, the policy denoises an action trajectory conditioned on recent observations. This is useful for manipulation because demonstrations often contain multimodal behavior: the same visual state may admit several valid reaching or grasping motions. The diffusion formulation also fits receding-horizon control, since the model can generate a short action sequence and execute only the first part before replanning. In the standard formulation, the denoising network receives an observation history $O_t$ through a visual encoder and learns a conditional action distribution $p_\theta(A_t \mid O_t)$.

Our work keeps this action-diffusion view intact. The difference is in the conditioning signal. Rather than conditioning the denoising process only on the current visual observation, we introduce a learned dynamics representation between observation and action generation. Given $O_t$, a recurrent world model infers a latent belief state. During denoising, candidate action sequences can also be rolled forward through the world model to produce imagined future latent states. The diffusion policy is therefore conditioned on both the present belief and the predicted consequences of candidate actions. This keeps the multimodal action modeling benefit of Diffusion Policy while adding an explicit predictive pathway.

\subsection{World model}

World models learn compact latent dynamics that summarize an agent's interaction with the environment. Earlier model-based methods often focused on low-dimensional states or proprioceptive observations~\cite{embed,chua}. For visual control, the model must learn a representation that is predictive enough for planning or policy learning while avoiding the cost of directly predicting pixels at every future step. Recurrent state-space models address this problem by combining a deterministic recurrent state with stochastic latent variables. The deterministic state carries information from the past, while the stochastic state captures uncertainty in the current latent transition.

Dreamer-style agents are a representative example of this approach. Dreamer learns a recurrent state-space model from image observations and trains behavior in the learned latent space through imagined rollouts~\cite{dreamer}. DreamerV3 further improves robustness across diverse domains using discrete stochastic latents, KL balancing, free bits, symlog transformations, two-hot reward and value prediction, and continuation modeling~\cite{dreamerv3}. The important idea for our setting is not only that the world model predicts future observations, but that it provides a latent space in which future action consequences can be evaluated before executing the actions.

We use this idea to augment a diffusion policy rather than replacing the policy with a standard actor-critic agent. The world model takes the observation history $O_t$ and infers a posterior latent state $s_t$. For a candidate action sequence $A_{t:t+H}$, the recurrent dynamics produce imagined latent states $\hat{s}_{t+1:t+H}$. These states are then used as additional conditioning for the diffusion policy:
\[
    p_\theta(A_t \mid O_t, s_t, \hat{s}_{t+1:t+H}).
\]
This differs from a pure world-model agent, where an actor is usually trained directly in latent imagination. It also differs from a standard diffusion policy, where the policy is reactive to the encoded observation history. Our design uses the world model as a predictive conditioning module: the diffusion model remains the action generator, while the world model provides the latent belief and imagined future context.

The resulting ablations separate the contribution of each component. A diffusion-only policy tests the original action diffusion formulation. An RSSM-only behavior cloning policy tests whether the learned world-model representation alone is sufficient for action prediction. A DP+RSSM policy tests whether the RSSM belief state improves diffusion-based action generation. Finally, dream training uses the RSSM's imagined rollouts and reward/value predictions to further guide the diffusion policy toward action candidates with better predicted futures.

% Suggested bibliography entries:
%
% \bibitem{ddpm}
% J. Ho, A. Jain, and P. Abbeel, ``Denoising diffusion probabilistic models,'' in \textit{Advances in Neural Information Processing Systems}, 2020.
%
% \bibitem{score}
% Y. Song, J. Sohl-Dickstein, D. P. Kingma, A. Kumar, S. Ermon, and B. Poole, ``Score-based generative modeling through stochastic differential equations,'' in \textit{International Conference on Learning Representations}, 2021.
%
% \bibitem{dp}
% C. Chi et al., ``Diffusion Policy: Visuomotor Policy Learning via Action Diffusion,'' \textit{arXiv preprint arXiv:2303.04137}, 2024.
%
% \bibitem{embed}
% M. Watter, J. Springenberg, J. Boedecker, and M. Riedmiller, ``Embed to control: A locally linear latent dynamics model for control from raw images,'' in \textit{Advances in Neural Information Processing Systems}, 2015.
%
% \bibitem{chua}
% K. Chua, R. Calandra, R. McAllister, and S. Levine, ``Deep reinforcement learning in a handful of trials using probabilistic dynamics models,'' in \textit{Advances in Neural Information Processing Systems}, 2018.
%
% \bibitem{dreamer}
% D. Hafner, T. Lillicrap, J. Ba, and M. Norouzi, ``Dream to Control: Learning Behaviors by Latent Imagination,'' in \textit{International Conference on Learning Representations}, 2020.
%
% \bibitem{dreamerv3}
% D. Hafner, J. Pasukonis, J. Ba, and T. Lillicrap, ``Mastering diverse domains through world models,'' \textit{arXiv preprint arXiv:2301.04104}, 2024.
